% GENERATED FILE. Edit canonical lessons, then run npm run generate:books.
% canonical-source-hash: fnv1a64:cd41e4f79bd947c8
% canonical-lessons: JA-C417-tenimotsu, JA-C417-nifuda, JA-C417-seou, JA-C417-kakaeru, JA-C417-korogasu

\chapter{Carrying the Luggage}
\label{ch:ja-carrying-the-luggage}
\hlchaptermodality{417}

\begin{chapteropening}
\textbf{By the end of this chapter:} \emph{I can carry, roll and label my luggage.}

Complete the last lesson of chapter 417: I can carry, roll and label my luggage.
\end{chapteropening}

\section[\texorpdfstring{\ja{てにもつ} (tenimotsu)}{てにもつ (tenimotsu)}]{\ja{てにもつ} (tenimotsu) — hand luggage}

\label{lesson:JA-C417-tenimotsu}



\begin{quote}
Before the new one: say the Japanese for to turn very cold, then the Japanese for to shine on.
\end{quote}

\subsection*{You'll want to know: \ja{てにもつ}}
\textbf{\ja{てにもつ}} — \emph{tenimotsu} — \textquotedblleft{}hand luggage\textquotedblright{}.

\begin{quote}
\textbf{\ja{てにもつを} \ja{ひとつ} \ja{もって、ひこうきに} \ja{のりました。}} — \emph{tenimotsu o hitotsu motte, hikōki ni norimashita} — \textquotedblleft{}I got on the plane with one piece of hand luggage.\textquotedblright{}
\end{quote}

\begin{cousinweb}
\textbf{\ja{てにもつ}} (\emph{tenimotsu}) is built from both layers of the vocabulary: the native Japanese (wago) \emph{te}, \textquotedblleft{}hand\textquotedblright{}, inherited rather than borrowed; the native Japanese (wago) \emph{ni}, \textquotedblleft{}load\textquotedblright{}, inherited rather than borrowed; and the Sino-Japanese \emph{motsu}, \textquotedblleft{}thing\textquotedblright{} (Middle Chinese \emph{mjut}), an on'yomi borrowed from Chinese (JMdict; KANJIDIC2; Middle Chinese in Baxter's transcription).
\end{cousinweb}

\begin{grammarlens}[title={What you've built}]
One more word for getting around.
\end{grammarlens}

\subsection*{Your turn}
\begin{itemize}
\raggedright
  \item \emph{Say it:} \emph{tenimotsu}
  \item \emph{Say it:} \emph{tenimotsu}, once more
  \item \emph{Say it:} \emph{tenimotsu o hitotsu motte, hikōki ni norimashita} — I got on the plane with one piece of hand luggage.
  \item \emph{Recall:} say the Japanese for to turn very cold, then the Japanese for to shine on, then say \emph{tenimotsu} again
\end{itemize}

\begin{tcolorbox}[breakable,colback=teal!4,colframe=teal!35!black,title={Before you move on}]
What does \ja{てにもつ} mean? (\textquotedblleft{}Hand luggage\textquotedblright{}.) Say it once more.
\end{tcolorbox}
\section[\texorpdfstring{\ja{にふだ} (nifuda)}{にふだ (nifuda)}]{\ja{にふだ} (nifuda) — a luggage tag}

\label{lesson:JA-C417-nifuda}



\begin{quote}
Before the new one: say the Japanese for to shine on, then the Japanese for hand luggage.
\end{quote}

\subsection*{You'll want to know: \ja{にふだ}}
\textbf{\ja{にふだ}} — \emph{nifuda} — \textquotedblleft{}a luggage tag\textquotedblright{}.

\begin{quote}
\textbf{\ja{にふだに} \ja{なまえを} \ja{かいて} \ja{ください。}} — \emph{nifuda ni namae o kaite kudasai} — \textquotedblleft{}Please write your name on the luggage tag.\textquotedblright{}
\end{quote}

\begin{cousinweb}
\textbf{\ja{にふだ}} (\emph{nifuda}) is a native Japanese word (wago), inherited rather than borrowed from Chinese, built from \emph{ni}, \textquotedblleft{}load\textquotedblright{} and \emph{fuda}, \textquotedblleft{}tag\textquotedblright{}; its kanji were attached later as kun'yomi (JMdict).
\end{cousinweb}

\begin{grammarlens}[title={What you've built}]
One more word for getting around.
\end{grammarlens}

\subsection*{Your turn}
\begin{itemize}
\raggedright
  \item \emph{Say it:} \emph{nifuda}
  \item \emph{Say it:} \emph{nifuda}, once more
  \item \emph{Say it:} \emph{nifuda ni namae o kaite kudasai} — Please write your name on the luggage tag.
  \item \emph{Recall:} say the Japanese for to shine on, then the Japanese for hand luggage, then say \emph{nifuda} again
\end{itemize}

\begin{tcolorbox}[breakable,colback=teal!4,colframe=teal!35!black,title={Before you move on}]
What does \ja{にふだ} mean? (\textquotedblleft{}A luggage tag\textquotedblright{}.) Say it once more.
\end{tcolorbox}
\section[\texorpdfstring{\ja{せおう} (seou)}{せおう (seou)}]{\ja{せおう} (seou) — to carry on one's back}

\label{lesson:JA-C417-seou}



\begin{quote}
Before the new one: say the Japanese for hand luggage, then the Japanese for a luggage tag.
\end{quote}

\subsection*{You'll want to know: \ja{せおう}}
\textbf{\ja{せおう}} — \emph{seou} — \textquotedblleft{}to carry on one's back\textquotedblright{}.

\begin{quote}
\textbf{\ja{おもい} \ja{にもつを} \ja{せおって、やまに} \ja{のぼりました。}} — \emph{omoi nimotsu o seotte, yama ni noborimashita} — \textquotedblleft{}I climbed the mountain with a heavy load on my back.\textquotedblright{}
\end{quote}

Polite \textbf{\ja{せおいます}}, \emph{seoimasu}; a godan verb. Past (polite): \textbf{\ja{あにが} \ja{にもつを} \ja{せおいました。}}, \emph{ani ga nimotsu o seoimashita}.

\begin{cousinweb}
\textbf{\ja{せおう}} (\emph{seou}) is a native Japanese word (wago), inherited rather than borrowed from Chinese, built from \emph{se}, \textquotedblleft{}back\textquotedblright{}, and \emph{ou}, \textquotedblleft{}to carry on one's back\textquotedblright{}; its kanji were attached later as kun'yomi (JMdict).
\end{cousinweb}

\begin{grammarlens}[title={What you've built}]
One more word for getting around.
\end{grammarlens}

\subsection*{Your turn}
\begin{itemize}
\raggedright
  \item \emph{Say it:} \emph{seou}
  \item \emph{Say it:} \emph{seou}, once more
  \item \emph{Say it:} \emph{omoi nimotsu o seotte, yama ni noborimashita} — I climbed the mountain with a heavy load on my back.
  \item \emph{Recall:} say the Japanese for hand luggage, then the Japanese for a luggage tag, then say \emph{seou} again
\end{itemize}

\begin{tcolorbox}[breakable,colback=teal!4,colframe=teal!35!black,title={Before you move on}]
What does \ja{せおう} mean? (\textquotedblleft{}To carry on one's back\textquotedblright{}.) Say it once more.
\end{tcolorbox}
\section[\texorpdfstring{\ja{かかえる} (kakaeru)}{かかえる (kakaeru)}]{\ja{かかえる} (kakaeru) — to carry in one's arms}

\label{lesson:JA-C417-kakaeru}



\begin{quote}
Before the new one: say the Japanese for a luggage tag, then the Japanese for to carry on one's back.
\end{quote}

\subsection*{You'll want to know: \ja{かかえる}}
\textbf{\ja{かかえる}} — \emph{kakaeru} — \textquotedblleft{}to carry in one's arms\textquotedblright{}.

\begin{quote}
\textbf{\ja{はこを} \ja{かかえて、かいだんを} \ja{のぼりました。}} — \emph{hako o kakaete, kaidan o noborimashita} — \textquotedblleft{}I carried a box up the stairs in my arms.\textquotedblright{}
\end{quote}

Polite \textbf{\ja{かかえます}}, \emph{kakaemasu}; an ichidan verb. Past (polite): \textbf{\ja{ちちが} \ja{はこを} \ja{かかえました。}}, \emph{chichi ga hako o kakaemashita}.

\begin{cousinweb}
\textbf{\ja{かかえる}} (\emph{kakaeru}) is a native Japanese word (wago), inherited rather than borrowed from Chinese; its kanji was attached later as a kun'yomi (JMdict).
\end{cousinweb}

\begin{grammarlens}[title={What you've built}]
One more word for getting around.
\end{grammarlens}

\subsection*{Your turn}
\begin{itemize}
\raggedright
  \item \emph{Say it:} \emph{kakaeru}
  \item \emph{Say it:} \emph{kakaeru}, once more
  \item \emph{Say it:} \emph{hako o kakaete, kaidan o noborimashita} — I carried a box up the stairs in my arms.
  \item \emph{Recall:} say the Japanese for a luggage tag, then the Japanese for to carry on one's back, then say \emph{kakaeru} again
\end{itemize}

\begin{tcolorbox}[breakable,colback=teal!4,colframe=teal!35!black,title={Before you move on}]
What does \ja{かかえる} mean? (\textquotedblleft{}To carry in one's arms\textquotedblright{}.) Say it once more.
\end{tcolorbox}
\section[\texorpdfstring{\ja{ころがす} (korogasu)}{ころがす (korogasu)}]{\ja{ころがす} (korogasu) — to roll; to wheel}

\label{lesson:JA-C417-korogasu}



\begin{quote}
Before the new one: say the Japanese for to carry on one's back, then the Japanese for to carry in one's arms.
\end{quote}

\subsection*{You'll want to know: \ja{ころがす}}
\textbf{\ja{ころがす}} — \emph{korogasu} — \textquotedblleft{}to roll; to wheel\textquotedblright{}.

\begin{quote}
\textbf{\ja{かばんを} \ja{ころがして、えきまで} \ja{あるきました。}} — \emph{kaban o korogashite, eki made arukimashita} — \textquotedblleft{}I wheeled my case to the station.\textquotedblright{}
\end{quote}

Polite \textbf{\ja{ころがします}}, \emph{korogashimasu}; a godan verb. Past (polite): \textbf{\ja{あねが} \ja{かばんを} \ja{ころがしました。}}, \emph{ane ga kaban o korogashimashita}.

\begin{cousinweb}
\textbf{\ja{ころがす}} (\emph{korogasu}) is a native Japanese word (wago), inherited rather than borrowed from Chinese; its kanji was attached later as a kun'yomi (JMdict).
\end{cousinweb}

\begin{grammarlens}[title={What you've built}]
One more word for getting around.
\end{grammarlens}

\subsection*{Your turn}
\begin{itemize}
\raggedright
  \item \emph{Say it:} \emph{korogasu}
  \item \emph{Say it:} \emph{korogasu}, once more
  \item \emph{Say it:} \emph{kaban o korogashite, eki made arukimashita} — I wheeled my case to the station.
  \item \emph{Recall:} say the Japanese for to carry on one's back, then the Japanese for to carry in one's arms, then say \emph{korogasu} again
\end{itemize}

\begin{tcolorbox}[breakable,colback=teal!4,colframe=teal!35!black,title={Before you move on}]
What does \ja{ころがす} mean? (\textquotedblleft{}To roll; to wheel\textquotedblright{}.) Say it once more.
\end{tcolorbox}
